% GENERATED FILE. Edit canonical lessons, then run npm run generate:books.
% canonical-source-hash: fnv1a64:d756cf51f9584606
% canonical-lessons: GU-C411-rajai, GU-C411-osiku, GU-C411-macchardani, GU-C411-sankelvu, GU-C411-suvdavvu

\chapter{Quilts and Pillows}
\label{ch:gu-quilts-and-pillows}
\hlchaptermodality{411}

\begin{chapteropening}
\textbf{By the end of this chapter:} \emph{I can ask for bedding and settle down for the night.}

Complete the last lesson of chapter 411: I can ask for bedding and settle down for the night.
\end{chapteropening}

\section[\texorpdfstring{\gu{રજાઈ} (rajāī)}{રજાઈ (rajāī)}]{\gu{રજાઈ} (rajāī) — a quilt}

\label{lesson:GU-C411-rajai}



\begin{quote}
Before the new one: say the Gujarati for to bail out, then the Gujarati for to wring out.
\end{quote}

\subsection*{You'll want to know: \gu{રજાઈ}}
\textbf{\gu{રજાઈ}} — \emph{rajāī} — \textquotedblleft{}a quilt\textquotedblright{}.

\begin{quote}
\gu{શિયાળામાં} \gu{રાત્રે} \gu{હું} \gu{ભારે} \gu{રજાઈ} \gu{ઓઢું} \gu{છું}.
\end{quote}

\begin{quote}
\emph{śiyāḷāmā̃ rātre hũ bhāre rajāī oḍhũ chũ.}
\end{quote}

\begin{quote}
On winter nights I wrap myself in a heavy quilt.
\end{quote}

A feminine noun.

\begin{cousinweb}
\textbf{\gu{રજાઈ}} (\emph{rajāī}) is usually traced to Persian \emph{razāī}, \textquotedblleft{}a quilt\textquotedblright{}, a cover stuffed with cotton.
\end{cousinweb}

\begin{grammarlens}[title={What you've built}]
One more word for getting around.
\end{grammarlens}

\subsection*{Your turn}
\begin{itemize}
\raggedright
  \item \emph{Say it:} \emph{rajāī}
  \item \emph{Say it:} \emph{rajāī}, once more
  \item \emph{Say it:} \emph{śiyāḷāmā̃ rātre hũ bhāre rajāī oḍhũ chũ.} — On winter nights I wrap myself in a heavy quilt.
  \item \emph{Recall:} say the Gujarati for to bail out, then the Gujarati for to wring out, then say \emph{rajāī} again
\end{itemize}

\begin{tcolorbox}[breakable,colback=teal!4,colframe=teal!35!black,title={Before you move on}]
What does \gu{રજાઈ} mean? (\textquotedblleft{}A quilt\textquotedblright{}.) Say it once more.
\end{tcolorbox}
\section[\texorpdfstring{\gu{ઓશીકું} (ośīkũ)}{ઓશીકું (ośīkũ)}]{\gu{ઓશીકું} (ośīkũ) — a pillow; a cushion for the head}

\label{lesson:GU-C411-osiku}



\begin{quote}
Before the new one: say the Gujarati for to wring out, then the Gujarati for a quilt.
\end{quote}

\subsection*{You'll want to know: \gu{ઓશીકું}}
\textbf{\gu{ઓશીકું}} — \emph{ośīkũ} — \textquotedblleft{}a pillow; a cushion for the head\textquotedblright{}.

\begin{quote}
\gu{મને} \gu{નરમ} \gu{ઓશીકું} \gu{ગમે} \gu{છે}.
\end{quote}

\begin{quote}
\emph{mane naram ośīkũ game che.}
\end{quote}

\begin{quote}
I like a soft pillow.
\end{quote}

A neuter noun.

\begin{cousinweb}
\textbf{\gu{ઓશીકું}} (\emph{ośīkũ}) goes back to Sanskrit \emph{ucchīrṣaka}, \textquotedblleft{}a pillow\textquotedblright{}, from a form for raising the head; CDIAL lists the Gujarati word (CDIAL 1853).
\end{cousinweb}

\begin{grammarlens}[title={What you've built}]
One more word for getting around.
\end{grammarlens}

\subsection*{Your turn}
\begin{itemize}
\raggedright
  \item \emph{Say it:} \emph{ośīkũ}
  \item \emph{Say it:} \emph{ośīkũ}, once more
  \item \emph{Say it:} \emph{mane naram ośīkũ game che.} — I like a soft pillow.
  \item \emph{Recall:} say the Gujarati for to wring out, then the Gujarati for a quilt, then say \emph{ośīkũ} again
\end{itemize}

\begin{tcolorbox}[breakable,colback=teal!4,colframe=teal!35!black,title={Before you move on}]
What does \gu{ઓશીકું} mean? (\textquotedblleft{}A pillow; a cushion for the head\textquotedblright{}.) Say it once more.
\end{tcolorbox}
\section[\texorpdfstring{\gu{મચ્છરદાની} (macchardānī)}{મચ્છરદાની (macchardānī)}]{\gu{મચ્છરદાની} (macchardānī) — a mosquito net}

\label{lesson:GU-C411-macchardani}



\begin{quote}
Before the new one: say the Gujarati for a quilt, then the Gujarati for a cushion for the head.
\end{quote}

\subsection*{You'll want to know: \gu{મચ્છરદાની}}
\textbf{\gu{મચ્છરદાની}} — \emph{macchardānī} — \textquotedblleft{}a mosquito net\textquotedblright{}.

\begin{quote}
\gu{રાત્રે} \gu{અમે} \gu{પલંગ} \gu{પર} \gu{મચ્છરદાની} \gu{બાંધીએ} \gu{છીએ}.
\end{quote}

\begin{quote}
\emph{rātre ame palaṅg par macchardānī bā̃dhīe chīe.}
\end{quote}

\begin{quote}
At night we hang a mosquito net over the bed.
\end{quote}

A feminine noun.

\begin{cousinweb}
\textbf{\gu{મચ્છરદાની}} (\emph{macchardānī}) is built from \gu{મચ્છર}, a mosquito, and -\gu{દાની}, \textquotedblleft{}a holder, a container\textquotedblright{}, an ending usually traced to Persian \emph{-dān}: a holder that keeps the mosquitoes out.
\end{cousinweb}

\begin{grammarlens}[title={What you've built}]
One more word for getting around.
\end{grammarlens}

\subsection*{Your turn}
\begin{itemize}
\raggedright
  \item \emph{Say it:} \emph{macchardānī}
  \item \emph{Say it:} \emph{macchardānī}, once more
  \item \emph{Say it:} \emph{rātre ame palaṅg par macchardānī bā̃dhīe chīe.} — At night we hang a mosquito net over the bed.
  \item \emph{Recall:} say the Gujarati for a quilt, then the Gujarati for a cushion for the head, then say \emph{macchardānī} again
\end{itemize}

\begin{tcolorbox}[breakable,colback=teal!4,colframe=teal!35!black,title={Before you move on}]
What does \gu{મચ્છરદાની} mean? (\textquotedblleft{}A mosquito net\textquotedblright{}.) Say it once more.
\end{tcolorbox}
\section[\texorpdfstring{\gu{સંકેલવું} (sankelvũ)}{સંકેલવું (sankelvũ)}]{\gu{સંકેલવું} (sankelvũ) — to roll up; to fold up (bedding)}

\label{lesson:GU-C411-sankelvu}



\begin{quote}
Before the new one: say the Gujarati for a cushion for the head, then the Gujarati for a mosquito net.
\end{quote}

\subsection*{You'll want to know: \gu{સંકેલવું}}
\textbf{\gu{સંકેલવું}} — \emph{sankelvũ} — \textquotedblleft{}to roll up; to fold up (bedding)\textquotedblright{}.

\begin{quote}
\gu{સવારે} \gu{ઉઠીને} \gu{મેં} \gu{પથારી} \gu{સંકેલી}.
\end{quote}

\begin{quote}
\emph{savāre uṭhīne mẽ pathārī sankelī.}
\end{quote}

\begin{quote}
In the morning I got up and rolled up the bedding.
\end{quote}

Past, with \gu{તેણે} (he or she …): \gu{તેણે} \gu{પથારી} \gu{સંકેલી}, \emph{teṇe pathārī sankelī}.

\begin{cousinweb}
\textbf{\gu{સંકેલવું}} (\emph{sankelvũ}) may belong with Hindi \emph{sakilnā}, \textquotedblleft{}to gather up, to shrink\textquotedblright{}, which CDIAL lists under Sanskrit \emph{saṁkirati}, \textquotedblleft{}mixes together\textquotedblright{} (CDIAL 12823); it gives no Gujarati form, so the link is hedged.
\end{cousinweb}

\begin{grammarlens}[title={What you've built}]
One more word for getting around.
\end{grammarlens}

\subsection*{Your turn}
\begin{itemize}
\raggedright
  \item \emph{Say it:} \emph{sankelvũ}
  \item \emph{Say it:} \emph{sankelvũ}, once more
  \item \emph{Say it:} \emph{savāre uṭhīne mẽ pathārī sankelī.} — In the morning I got up and rolled up the bedding.
  \item \emph{Recall:} say the Gujarati for a cushion for the head, then the Gujarati for a mosquito net, then say \emph{sankelvũ} again
\end{itemize}

\begin{tcolorbox}[breakable,colback=teal!4,colframe=teal!35!black,title={Before you move on}]
What does \gu{સંકેલવું} mean? (\textquotedblleft{}To roll up; to fold up (bedding)\textquotedblright{}.) Say it once more.
\end{tcolorbox}
\section[\texorpdfstring{\gu{સુવડાવવું} (suvḍāvvũ)}{સુવડાવવું (suvḍāvvũ)}]{\gu{સુવડાવવું} (suvḍāvvũ) — to put to sleep (a child)}

\label{lesson:GU-C411-suvdavvu}



\begin{quote}
Before the new one: say the Gujarati for a mosquito net, then the Gujarati for to roll up.
\end{quote}

\subsection*{You'll want to know: \gu{સુવડાવવું}}
\textbf{\gu{સુવડાવવું}} — \emph{suvḍāvvũ} — \textquotedblleft{}to put to sleep (a child)\textquotedblright{}.

\begin{quote}
\gu{માતા} \gu{બાળકને} \gu{ગીત} \gu{ગાઈને} \gu{સુવડાવે} \gu{છે}.
\end{quote}

\begin{quote}
\emph{mātā bāḷakne gīt gāīne suvḍāve che.}
\end{quote}

\begin{quote}
The mother sings the child to sleep.
\end{quote}

Past, with \gu{તેણે} (he or she …): \gu{તેણે} \gu{બાળકને} \gu{સુવડાવ્યું}, \emph{teṇe bāḷakne suvḍāvyũ}.

\begin{cousinweb}
\textbf{\gu{સુવડાવવું}} (\emph{suvḍāvvũ}) is the causative partner of \gu{સૂવું}, \textquotedblleft{}to sleep, to lie down\textquotedblright{}, a tadbhava of Sanskrit \emph{svápati}, \textquotedblleft{}sleeps\textquotedblright{} (CDIAL 13902): to make someone sleep.
\end{cousinweb}

\begin{grammarlens}[title={What you've built}]
One more word for getting around.
\end{grammarlens}

\subsection*{Your turn}
\begin{itemize}
\raggedright
  \item \emph{Say it:} \emph{suvḍāvvũ}
  \item \emph{Say it:} \emph{suvḍāvvũ}, once more
  \item \emph{Say it:} \emph{mātā bāḷakne gīt gāīne suvḍāve che.} — The mother sings the child to sleep.
  \item \emph{Recall:} say the Gujarati for a mosquito net, then the Gujarati for to roll up, then say \emph{suvḍāvvũ} again
\end{itemize}

\begin{tcolorbox}[breakable,colback=teal!4,colframe=teal!35!black,title={Before you move on}]
What does \gu{સુવડાવવું} mean? (\textquotedblleft{}To put to sleep (a child)\textquotedblright{}.) Say it once more.
\end{tcolorbox}
